\documentclass[11pt]{article}
\usepackage[margin=1in]{geometry}
\usepackage{amsmath,amssymb,amsthm,hyperref}
\newtheorem{theorem}{Theorem}
\newtheorem{corollary}{Corollary}
\title{Nonintegrality of Schmidt-type coefficients}
\author{Principia Math In Progress\\\small AI-assisted research note; author review only}
\date{27 September 2026}
\begin{document}
\maketitle
\begin{abstract}
We prove Guo and Zeng's nonintegrality conjecture for every pair of integers
$r>s>1$. The defining equation at $n=p-2$, reduced at a prime
$p\ge5$ not dividing $r-s$, gives the obstruction. In fact the least common
multiple of the first $p-1$ coefficient denominators is divisible by
$p^{s-1}$. Thus infinitely many coefficients are nonintegral, and all but
finitely many primes occur in their denominators.
\end{abstract}

\section*{Task 1: the original conjecture}
For positive integers $r,s$, let $c_k=c_k^{(r,s)}$ be the unique rational
numbers defined by
\begin{equation}\label{eq:definition}
 \sum_{k=0}^n B_{n,k}^{\,r}
 =\sum_{k=0}^n B_{n,k}^{\,s}c_k\quad(n\ge0),\qquad
 B_{n,k}=\binom nk\binom{n+k}{k}.
\end{equation}
The diagonal is nonzero, so this is a triangular definition.
Conjecture 4.2 of Guo and Zeng~\cite{GZ}, recorded as
\href{https://mathdb.com/p/323297}{MathDB 323297}, states that for every
$r>s>1$ there exists a positive index $k$ for which $c_k$ is not an integer.
We prove the following stronger assertion.

\begin{theorem}\label{thm:denom}
Fix integers $r>s>1$ and put $d=r-s$. For each prime $p\ge5$ with $p\nmid d$,
let $D_{p-2}$ be the least common multiple of the positive reduced
denominators of $c_0,\ldots,c_{p-2}$. Then
\[
 p^{s-1}\mid D_{p-2}.
\]
In particular some $c_k$, with $2\le k\le p-2$, has reduced denominator
divisible by $p^{s-1}$.
\end{theorem}

\begin{proof}
The cases $n=0,1$ of~\eqref{eq:definition} give
\begin{equation}\label{eq:first}
 c_0=1,\qquad c_1=2^{r-s}=2^d.
\end{equation}
Set $n=p-2$ and $a=B_{p-2,1}=(p-2)(p-1)=2-3p+p^2$.
For every $2\le k\le p-2$, the binomial coefficient
\[
 \binom{p-2+k}{k}=\frac{(p-1)p(p+1)\cdots(p-2+k)}{k!}
\]
contains a factor $p$ in the numerator and none in the denominator.
Consequently $p\mid B_{p-2,k}$. (The numerator has exactly $k$ factors;
the displayed product begins at $p-1$.)

Write $D=D_{p-2}$. After using~\eqref{eq:first}, multiplying the equation
at $n=p-2$ by $D$, and separating $k=0,1$, we obtain the integer identity
\begin{equation}\label{eq:cleared}
 D a^s(a^d-2^d)
 +D\sum_{k=2}^{p-2}B_{p-2,k}^{\,r}
 =\sum_{k=2}^{p-2}B_{p-2,k}^{\,s}(Dc_k).
\end{equation}
Every $Dc_k$ is an integer. The right side is divisible by $p^s$.
The second term on the left is divisible by $p^r$, and $r>s$.
It follows that
\begin{equation}\label{eq:divisibility}
 p^s\mid D a^s(a^d-2^d).
\end{equation}
Since $a\equiv2\pmod p$, $p\nmid a$. The binomial theorem gives
\[
 a^d-2^d\equiv(2-3p)^d-2^d
 \equiv-3d\,2^{d-1}p\pmod{p^2}.
\]
The coefficient $-3d2^{d-1}$ is nonzero modulo $p$, because $p\ge5$ and
$p\nmid d$. Thus $a^d-2^d$ is divisible by $p$ but not by $p^2$.
The divisibility~\eqref{eq:divisibility} therefore forces $p^{s-1}\mid D$.
An lcm attains each prime exponent at one of its factors, which proves
the last assertion. The indices 0 and 1 have denominator 1 by~\eqref{eq:first}.
\end{proof}

\begin{corollary}\label{cor:conjecture}
Guo and Zeng's Conjecture 4.2 holds. More strongly, for each fixed pair
$r>s>1$ there are infinitely many indices $k$ with $c_k^{(r,s)}\notin\mathbb Z$.
Every prime $p\ge5$ not dividing $r-s$ divides the reduced denominator of
some coefficient, to at least the power $s-1$.
\end{corollary}
\begin{proof}
There are infinitely many primes not dividing the fixed nonzero integer
$6(r-s)$. Apply Theorem~\ref{thm:denom} to each. A finite collection of
rational numbers has only finitely many prime factors in its reduced
denominators. Hence infinitely many distinct nonintegral coefficients
are necessary. In particular at least one exists, as conjectured.
\end{proof}

An explicit upper bound for the first nonintegral index is $p_0-2$, where
$p_0$ is the least prime at least 5 not dividing $r-s$. Also, writing $D_N$
for the denominator lcm through index $N$, we immediately obtain
\[
 \prod_{\substack{5\le p\le N+2\\p\text{ prime},\ p\nmid(r-s)}}p^{s-1}
 \ \mid\ D_N.
\]
These claims follow from the same proof and receive no separate problem credit.

\subsection*{Scope, source and checks}
The full current author manuscript and the full final published paper
were read; the definition and Conjecture 4.2 agree. The different finite
integrality conjecture, Conjecture 4.1, has its own proof package and is
not counted again here. The theorem above is a direct proof for all
integers $r>s>1$; it does not depend on that other result or on computation.

The accompanying \texttt{verify\_denominators.py} performs supplementary
exact checks of the binomial divisibility, first-order obstruction, and
actual rational coefficient denominators over stated finite ranges.
These checks test conventions and edge cases; the universal conclusion
is proved above. The literature audit records a dated, bounded search
for prior claims. No independent peer review or formal proof-assistant
verification is claimed.

\section*{Four further research tasks}
\paragraph{Task 2.} For each fixed $r>s>1$, determine the exact largest
$p$-power in $D_{p-2}$ for every $p\ge5$, including primes dividing $r-s$.
\paragraph{Task 3.} Determine the asymptotic growth of $\log D_N$ for each
fixed pair $r>s>1$, sharpening the prime-product lower bound above.
\paragraph{Task 4.} Use the equations at $n=p-m$, with fixed $m\ge3$, to
derive additional necessary congruences on $r,s$ for integrality through
$p-m$, and compare their strength with the $m=2$ obstruction.
\paragraph{Task 5.} Determine the first nonintegral index as a function
of $(r,s)$ and characterize the pairs for which the bound $p_0-2$ is attained.

\begin{thebibliography}{9}
\bibitem{GZ} V. J. W. Guo and J. Zeng,
\emph{On Zudilin's $q$-question about Schmidt's problem},
Electronic Journal of Combinatorics \textbf{19}(3) (2012), P4,
\href{https://doi.org/10.37236/2231}{doi:10.37236/2231}.
Conjecture 4.2 appears on page 4 of the final article;
current author manuscript:
\href{https://arxiv.org/abs/1204.0187v3}{arXiv:1204.0187v3}.
\end{thebibliography}
\end{document}
